% ============================================================================
% EMCH 792: Optimal State Estimation - Homework 3 (Written Portion)
% ============================================================================

\documentclass[12pt]{article}

% ============================================================================
% PACKAGE IMPORTS
% ============================================================================
\usepackage{amsmath}
\usepackage{amssymb}
\usepackage[margin=1in]{geometry}
\usepackage{setspace}
\usepackage{xcolor}
\usepackage{enumitem}
\usepackage{tcolorbox}
\usepackage{fancyhdr}

% ============================================================================
% HEADER AND FOOTER SETUP
% ============================================================================
\pagestyle{fancy}
\setlength{\headheight}{14.5pt}
\fancyhf{}
\fancyhead[L]{Page \thepage}
\fancyhead[C]{EMCH 792}
\fancyhead[R]{J.C. Vaught}
\renewcommand{\headrulewidth}{0pt}

% --- USC Brand Colors ---
\definecolor{USC_Garnet}{HTML}{73000A}
\definecolor{USC_Black90}{HTML}{363636}
\definecolor{USC_Black10}{HTML}{ECECEC}
\definecolor{USC_Honeycomb}{HTML}{A49137}
\definecolor{USC_Horseshoe}{HTML}{65780B}

% ============================================================================
% CUSTOM BOXES
% ============================================================================
\newtcolorbox{stepbox}{
  colback=white,
  colframe=USC_Garnet,
  fonttitle=\bfseries,
  title=Derivation,
  sharp corners,
  colbacktitle=USC_Garnet,
  coltitle=white,
}

\newtcolorbox{resultsbox}{
  colback=white,
  colframe=USC_Horseshoe,
  fonttitle=\bfseries,
  title=Final Result,
  sharp corners,
  colbacktitle=USC_Horseshoe,
  coltitle=white,
}

\newcommand{\question}[1]{%
  \clearpage
  \vspace{0.5cm}
  {\noindent\normalsize \textbf{#1}}
  \vspace{0.2cm}
  \hrule
  \vspace{0.1cm}
  \hrule
  \vspace{0.3cm}
}

% ============================================================================
% DOCUMENT INFORMATION
% ============================================================================
\title{EMCH 792: Optimal State Estimation \\ Homework 3 \\ \large{Written Portion Solutions}}
\author{Instructor: N. Vitzilaios, Department of Mechanical Engineering \\ University of South Carolina \\ \textit{Solutions by: J.C. Vaught}}
\date{Due: February 26, 2026}

\begin{document}

\maketitle
\setlength{\parindent}{0pt}
\setlist[enumerate,1]{label=\arabic*.}
\setlist[enumerate,2]{label=(\alph*)}

\begin{center}
\begin{tcolorbox}[colback=white, colframe=gray!50!black, colbacktitle=gray!20!white,
                  coltitle=black, sharp corners, boxrule=1pt,
                  title=\Large\bfseries Table of Contents]
\begin{tabular}{p{0.82\textwidth}r}
\textbf{Exercise 2.5:} Exponential Random Variable (20 pts) \dotfill & \textbf{\pageref{ex:2.5}} \\[0.3em]
\textbf{Exercise 2.8:} Sum of Uncorrelated RVs (20 pts) \dotfill & \textbf{\pageref{ex:2.8}} \\[0.3em]
\textbf{Exercise 2.12:} Process $Y(t)=X\cos t$ (20 pts) \dotfill & \textbf{\pageref{ex:2.12}} \\
\end{tabular}
\vspace{0.2cm}
\end{tcolorbox}
\end{center}

\question{Exercise 2.5 (Page 75): Exponentially Distributed Random Variable}\label{ex:2.5}
Given
\[
f_X(x)=
\begin{cases}
\alpha e^{-\alpha x}, & x>0,\\
0, & x\le 0,
\end{cases}
\qquad \alpha > 0.
\]

\begin{stepbox}
\textbf{(a) Find the CDF $F_X(x)$}

By definition,
\[
F_X(x)=P(X\le x)=\int_{-\infty}^{x}f_X(z)\,dz.
\]
Because the density is zero for $z\le0$, evaluate in two regions:
\[
F_X(x)=0,\qquad x\le0.
\]
For $x>0$,
\[
F_X(x)=\int_0^x \alpha e^{-\alpha z}\,dz
=\left[-e^{-\alpha z}\right]_{0}^{x}
=1-e^{-\alpha x}.
\]
Therefore,
\[
F_X(x)=
\begin{cases}
0, & x\le0,\\
1-e^{-\alpha x}, & x>0.
\end{cases}
\]
\end{stepbox}

\begin{stepbox}
\textbf{(b) Find the mean $\mu_X=E[X]$}

Start with
\[
E[X]=\int_0^{\infty}x\alpha e^{-\alpha x}\,dx.
\]
Use integration by parts with
\[
u=x,\quad dv=\alpha e^{-\alpha x}dx,
\quad du=dx,\quad v=-e^{-\alpha x}.
\]
Then
\[
E[X]=\left[-xe^{-\alpha x}\right]_0^{\infty}+\int_0^{\infty}e^{-\alpha x}\,dx.
\]
The boundary term is zero, and
\[
\int_0^{\infty}e^{-\alpha x}\,dx
=\left[-\frac{1}{\alpha}e^{-\alpha x}\right]_0^{\infty}
=\frac{1}{\alpha}.
\]
So
\[
\mu_X=E[X]=\frac{1}{\alpha}.
\]
\end{stepbox}

\begin{stepbox}
\textbf{(c) Find the second moment $E[X^2]$}

\[
E[X^2]=\int_0^{\infty}x^2\alpha e^{-\alpha x}\,dx.
\]
Integrate by parts once:
\[
u=x^2,\ dv=\alpha e^{-\alpha x}dx
\Rightarrow du=2x\,dx,\ v=-e^{-\alpha x}.
\]
Hence
\[
E[X^2]=\left[-x^2e^{-\alpha x}\right]_0^{\infty}+2\int_0^{\infty}xe^{-\alpha x}\,dx.
\]
Define
\[
I=\int_0^{\infty}xe^{-\alpha x}\,dx.
\]
Integrate $I$ by parts:
\[
u=x,\ dv=e^{-\alpha x}dx
\Rightarrow du=dx,\ v=-\frac{1}{\alpha}e^{-\alpha x}.
\]
So
\[
I=\left[-\frac{x}{\alpha}e^{-\alpha x}\right]_0^{\infty}+\frac{1}{\alpha}\int_0^{\infty}e^{-\alpha x}dx
=0+\frac{1}{\alpha}\cdot\frac{1}{\alpha}
=\frac{1}{\alpha^2}.
\]
Therefore,
\[
E[X^2]=2I=\frac{2}{\alpha^2}.
\]
\end{stepbox}

\begin{stepbox}
\textbf{(d) Find the variance $\sigma_X^2$}

Use
\[
\sigma_X^2=E[X^2]-(E[X])^2.
\]
Substitute parts (b) and (c):
\[
\sigma_X^2=\frac{2}{\alpha^2}-\left(\frac{1}{\alpha}\right)^2
=\frac{1}{\alpha^2}.
\]
Thus,
\[
\sigma_X=\sqrt{\sigma_X^2}=\frac{1}{\alpha}.
\]
\end{stepbox}

\begin{stepbox}
\textbf{(e) Probability that $X$ is within one standard deviation of the mean}

From above,
\[
\mu_X=\frac{1}{\alpha},\qquad \sigma_X=\frac{1}{\alpha}.
\]
So
\[
\mu_X-\sigma_X=0,
\qquad
\mu_X+\sigma_X=\frac{2}{\alpha}.
\]
Hence
\[
P(\mu_X-\sigma_X\le X\le \mu_X+\sigma_X)
=P\left(0\le X\le\frac{2}{\alpha}\right)
=F_X\left(\frac{2}{\alpha}\right)-F_X(0).
\]
Using the CDF,
\[
F_X\left(\frac{2}{\alpha}\right)=1-e^{-2},
\qquad F_X(0)=0.
\]
Therefore
\[
P\left(|X-\mu_X|\le\sigma_X\right)=1-e^{-2}\approx0.8646647.
\]
\end{stepbox}

\begin{resultsbox}
\[
F_X(x)=\begin{cases}0,&x\le0\\1-e^{-\alpha x},&x>0\end{cases},
\quad
E[X]=\frac{1}{\alpha},
\quad
E[X^2]=\frac{2}{\alpha^2},
\quad
\mathrm{Var}(X)=\frac{1}{\alpha^2},
\]
\[
P\left(|X-\mu_X|\le\sigma_X\right)=1-e^{-2}\approx0.8647.
\]
\end{resultsbox}

\question{Exercise 2.8 (Page 76): Sum of Two Zero-Mean Uncorrelated RVs}\label{ex:2.8}
Let $W$ and $V$ be zero-mean, uncorrelated RVs with standard deviations $\sigma_W$ and $\sigma_V$. Define
\[
X=W+V.
\]

\begin{stepbox}
Start from the definition of variance:
\[
\mathrm{Var}(X)=E[(X-E[X])^2].
\]
Because $E[W]=E[V]=0$,
\[
E[X]=E[W+V]=E[W]+E[V]=0.
\]
So
\[
\mathrm{Var}(X)=E[(W+V)^2]
=E[W^2+2WV+V^2]
=E[W^2]+2E[WV]+E[V^2].
\]
For uncorrelated variables,
\[
\mathrm{Cov}(W,V)=E[(W-E[W])(V-E[V])]=E[WV]=0.
\]
Hence
\[
\mathrm{Var}(X)=E[W^2]+E[V^2]
=\sigma_W^2+\sigma_V^2.
\]
Taking square root (standard deviation is nonnegative):
\[
\sigma_X=\sqrt{\sigma_W^2+\sigma_V^2}.
\]
\end{stepbox}

\begin{resultsbox}
\[
\boxed{\sigma_X = \sqrt{\sigma_W^2 + \sigma_V^2}}
\]
\end{resultsbox}

\question{Exercise 2.12 (Page 76): $Y(t)=X\cos t$}\label{ex:2.12}
Suppose $X$ is a random variable and
\[
Y(t)=X\cos t.
\]

\begin{stepbox}
\textbf{(a) Expected value $E[Y(t)]$}

For fixed $t$, $\cos t$ is deterministic, so it can be pulled out of expectation:
\[
E[Y(t)] = E[X\cos t] = \cos t\,E[X].
\]
If $\bar x := E[X]$, then
\[
E[Y(t)] = \bar x\cos t.
\]
\end{stepbox}

\begin{stepbox}
\textbf{(b) Time average $A[Y(t)]$}

For a particular realization $x$ of $X$, the signal is $y(t)=x\cos t$. Its time average is
\[
A[y(t)]
=\lim_{T\to\infty}\frac{1}{T}\int_0^T x\cos t\,dt
=x\lim_{T\to\infty}\frac{1}{T}\left[\sin t\right]_0^T
=x\lim_{T\to\infty}\frac{\sin T}{T}.
\]
Since $|\sin T|\le1$, we have $\sin T/T\to0$, therefore
\[
A[y(t)] = 0.
\]
In shorthand notation this is often written as
\[
A[Y(t)] = 0.
\]
\end{stepbox}

\begin{stepbox}
\textbf{(c) Condition under which $y(t)=A[Y(t)]$}

Since $A[Y(t)]=0$, the condition $y(t)=A[Y(t)]$ for all $t$ becomes
\[
x\cos t = 0 \quad \forall t.
\]
This is true for all $t$ only if
\[
x=0.
\]
If the question is interpreted as equality between ensemble mean and time average, then
\[
E[Y(t)] = A[Y(t)]
\iff \bar x\cos t = 0 \ \forall t
\iff \bar x=0.
\]
\end{stepbox}

\begin{resultsbox}
\[
E[Y(t)] = \bar x\cos t,
\qquad
A[Y(t)] = 0.
\]
For each realization, $y(t)=A[Y(t)]$ for all $t$ only when $x=0$. For mean/time-average equality for all $t$, the required condition is $E[X]=\bar x=0$.
\end{resultsbox}

\end{document}
